% ======================================================================
%  APÉNDICE C — GF(16) DE BOLSILLO
% ======================================================================

\chapter{GF(16) de bolsillo}\label{app:gf16}

El cuerpo $\FF4=\F_2[z]/(z^4+z+1)$ es el banco de pruebas de todo el
libro (ejemplo~\ref{ej:gf16} y los laboratorios de los
capítulos~\ref{ch:labgoppa} y~\ref{ch:labkem}). Este apéndice reúne,
para consulta rápida, las tablas que los ejercicios y las cuentas a
mano necesitan: potencias, logaritmos, inversos, raíces cuadradas,
traza y la tabla de multiplicar completa. Convenio: cada elemento se
escribe como dígito hexadecimal con el bit $i$-ésimo igual al
coeficiente de $z^i$ (en particular $z=\mathtt{2}$ y $z^4=z+1$); los
dígitos hex van en minúsculas, como en el programa del
apéndice~\ref{app:codigo}.

\section{Potencias y logaritmos}

\begin{table}[htbp]
\centering
\begin{tabular}{@{}clc@{\hspace{3.5em}}clc@{}}
\toprule
$i$ & hex & polinomio & $i$ & hex & polinomio\\
\midrule
0 & \texttt{1} & $1$ & 8 & \texttt{5} & $1+z^2$\\
1 & \texttt{2} & $z$ & 9 & \texttt{a} & $z+z^3$\\
2 & \texttt{4} & $z^2$ & 10 & \texttt{7} & $1+z+z^2$\\
3 & \texttt{8} & $z^3$ & 11 & \texttt{e} & $z+z^2+z^3$\\
4 & \texttt{3} & $1+z$ & 12 & \texttt{f} & $1+z+z^2+z^3$\\
5 & \texttt{6} & $z+z^2$ & 13 & \texttt{d} & $1+z^2+z^3$\\
6 & \texttt{c} & $z^2+z^3$ & 14 & \texttt{9} & $1+z^3$\\
7 & \texttt{b} & $1+z+z^3$ &  & &\\
\bottomrule
\end{tabular}
\caption{Las potencias $z^i$, $i=0,\dots,14$: el valor de $z^i$ es el
dígito hex y su expresión polinómica, la de al lado. El cuerpo tiene
$15$ elementos no nulos: $z$ es primitivo y su ciclo los agota todos
(figura~\ref{fig:ciclo-z}).}
\end{table}

\begin{table}[htbp]
\centering
\begin{tabular}{@{}cc@{\hspace{4em}}cc@{}}
\toprule
$a$ & $\log_z a$ & $a$ & $\log_z a$\\
\midrule
\texttt{1} & \texttt{0} & \texttt{9} & \texttt{e}\\
\texttt{2} & \texttt{1} & \texttt{a} & \texttt{9}\\
\texttt{3} & \texttt{4} & \texttt{b} & \texttt{7}\\
\texttt{4} & \texttt{2} & \texttt{c} & \texttt{6}\\
\texttt{5} & \texttt{8} & \texttt{d} & \texttt{d}\\
\texttt{6} & \texttt{5} & \texttt{e} & \texttt{b}\\
\texttt{7} & \texttt{a} & \texttt{f} & \texttt{c}\\
\texttt{8} & \texttt{3} & \texttt{0} & ---\\
\bottomrule
\end{tabular}
\caption{Logaritmos en base $z$ (los exponentes, también en hex: por
ejemplo $\log_z\mathtt{6}=\mathtt{5}$, pues $z^5=z+z^2$, y el $0$ no
tiene logaritmo). Para multiplicar: $ab=z^{\log a+\log b}$ con el
exponente módulo $15$.}
\end{table}

\section{Inversos, cuadrados, raíces y traza}

\begin{table}[htbp]
\centering
\begin{tabular}{@{}cc@{\hspace{4em}}cc@{}}
\toprule
$a$ & $a^{-1}$ & $a$ & $a^{-1}$\\
\midrule
\texttt{1} & \texttt{1} & \texttt{9} & \texttt{2}\\
\texttt{2} & \texttt{9} & \texttt{a} & \texttt{c}\\
\texttt{3} & \texttt{e} & \texttt{b} & \texttt{5}\\
\texttt{4} & \texttt{d} & \texttt{c} & \texttt{a}\\
\texttt{5} & \texttt{b} & \texttt{d} & \texttt{4}\\
\texttt{6} & \texttt{7} & \texttt{e} & \texttt{3}\\
\texttt{7} & \texttt{6} & \texttt{f} & \texttt{8}\\
\texttt{8} & \texttt{f} & \texttt{0} & ---\\
\bottomrule
\end{tabular}
\caption{Inversos multiplicativos. Comprobación rápida:
$\mathtt{3}\cdot\mathtt{e}=z^{4}\cdot z^{11}=z^{15}=1$.}
\end{table}

\begin{table}[htbp]
\centering
\begin{tabular}{@{}cccc@{\hspace{3em}}cccc@{}}
\toprule
$x$ & $x^2$ & $\sqrt{x}$ & $\Tr(x)$ & $x$ & $x^2$ & $\sqrt{x}$ & $\Tr(x)$\\
\midrule
\texttt{0} & \texttt{0} & \texttt{0} & 0 & \texttt{8} & \texttt{c} & \texttt{a} & 1\\
\texttt{1} & \texttt{1} & \texttt{1} & 0 & \texttt{9} & \texttt{d} & \texttt{b} & 1\\
\texttt{2} & \texttt{4} & \texttt{5} & 0 & \texttt{a} & \texttt{8} & \texttt{f} & 1\\
\texttt{3} & \texttt{5} & \texttt{4} & 0 & \texttt{b} & \texttt{9} & \texttt{e} & 1\\
\texttt{4} & \texttt{3} & \texttt{2} & 0 & \texttt{c} & \texttt{f} & \texttt{8} & 1\\
\texttt{5} & \texttt{2} & \texttt{3} & 0 & \texttt{d} & \texttt{e} & \texttt{9} & 1\\
\texttt{6} & \texttt{7} & \texttt{7} & 0 & \texttt{e} & \texttt{b} & \texttt{d} & 1\\
\texttt{7} & \texttt{6} & \texttt{6} & 0 & \texttt{f} & \texttt{a} & \texttt{c} & 1\\
\bottomrule
\end{tabular}
\caption{Cuadrados, raíces cuadradas ($\sqrt{x}=x^{8}$, la mitad de
$2^{m-1}$ porque la característica es $2$) y traza absoluta
($\Tr(x)=x+x^2+x^4+x^8\in\F_2$). Exactamente $8$ elementos tienen
traza $1$: el núcleo de $\Tr$ es un subespacio de dimensión $3$.}
\end{table}

\section{La tabla de multiplicar}

\begin{table}[htbp]
\centering
\setlength{\tabcolsep}{4.2pt}
\begin{tabular}{@{}c@{}c@{}c@{}c@{}c@{}c@{}c@{}c@{}c@{}c@{}c@{}c@{}c@{}c@{}c@{}c@{}c@{}}
\toprule
$\cdot$ & \texttt{0} & \texttt{1} & \texttt{2} & \texttt{3} & \texttt{4} & \texttt{5} & \texttt{6} & \texttt{7} & \texttt{8} & \texttt{9} & \texttt{a} & \texttt{b} & \texttt{c} & \texttt{d} & \texttt{e} & \texttt{f}\\
\midrule
\texttt{0} & \texttt{0} & \texttt{0} & \texttt{0} & \texttt{0} & \texttt{0} & \texttt{0} & \texttt{0} & \texttt{0} & \texttt{0} & \texttt{0} & \texttt{0} & \texttt{0} & \texttt{0} & \texttt{0} & \texttt{0} & \texttt{0}\\
\texttt{1} & \texttt{0} & \texttt{1} & \texttt{2} & \texttt{3} & \texttt{4} & \texttt{5} & \texttt{6} & \texttt{7} & \texttt{8} & \texttt{9} & \texttt{a} & \texttt{b} & \texttt{c} & \texttt{d} & \texttt{e} & \texttt{f}\\
\texttt{2} & \texttt{0} & \texttt{2} & \texttt{4} & \texttt{6} & \texttt{8} & \texttt{a} & \texttt{c} & \texttt{e} & \texttt{3} & \texttt{1} & \texttt{7} & \texttt{5} & \texttt{b} & \texttt{9} & \texttt{f} & \texttt{d}\\
\texttt{3} & \texttt{0} & \texttt{3} & \texttt{6} & \texttt{5} & \texttt{c} & \texttt{f} & \texttt{a} & \texttt{9} & \texttt{b} & \texttt{8} & \texttt{d} & \texttt{e} & \texttt{7} & \texttt{4} & \texttt{1} & \texttt{2}\\
\texttt{4} & \texttt{0} & \texttt{4} & \texttt{8} & \texttt{c} & \texttt{3} & \texttt{7} & \texttt{b} & \texttt{f} & \texttt{6} & \texttt{2} & \texttt{e} & \texttt{a} & \texttt{5} & \texttt{1} & \texttt{d} & \texttt{9}\\
\texttt{5} & \texttt{0} & \texttt{5} & \texttt{a} & \texttt{f} & \texttt{7} & \texttt{2} & \texttt{d} & \texttt{8} & \texttt{e} & \texttt{b} & \texttt{4} & \texttt{1} & \texttt{9} & \texttt{c} & \texttt{3} & \texttt{6}\\
\texttt{6} & \texttt{0} & \texttt{6} & \texttt{c} & \texttt{a} & \texttt{b} & \texttt{d} & \texttt{7} & \texttt{1} & \texttt{5} & \texttt{3} & \texttt{9} & \texttt{f} & \texttt{e} & \texttt{8} & \texttt{2} & \texttt{4}\\
\texttt{7} & \texttt{0} & \texttt{7} & \texttt{e} & \texttt{9} & \texttt{f} & \texttt{8} & \texttt{1} & \texttt{6} & \texttt{d} & \texttt{a} & \texttt{3} & \texttt{4} & \texttt{2} & \texttt{5} & \texttt{c} & \texttt{b}\\
\texttt{8} & \texttt{0} & \texttt{8} & \texttt{3} & \texttt{b} & \texttt{6} & \texttt{e} & \texttt{5} & \texttt{d} & \texttt{c} & \texttt{4} & \texttt{f} & \texttt{7} & \texttt{a} & \texttt{2} & \texttt{9} & \texttt{1}\\
\texttt{9} & \texttt{0} & \texttt{9} & \texttt{1} & \texttt{8} & \texttt{2} & \texttt{b} & \texttt{3} & \texttt{a} & \texttt{4} & \texttt{d} & \texttt{5} & \texttt{c} & \texttt{6} & \texttt{f} & \texttt{7} & \texttt{e}\\
\texttt{a} & \texttt{0} & \texttt{a} & \texttt{7} & \texttt{d} & \texttt{e} & \texttt{4} & \texttt{9} & \texttt{3} & \texttt{f} & \texttt{5} & \texttt{8} & \texttt{2} & \texttt{1} & \texttt{b} & \texttt{6} & \texttt{c}\\
\texttt{b} & \texttt{0} & \texttt{b} & \texttt{5} & \texttt{e} & \texttt{a} & \texttt{1} & \texttt{f} & \texttt{4} & \texttt{7} & \texttt{c} & \texttt{2} & \texttt{9} & \texttt{d} & \texttt{6} & \texttt{8} & \texttt{3}\\
\texttt{c} & \texttt{0} & \texttt{c} & \texttt{b} & \texttt{7} & \texttt{5} & \texttt{9} & \texttt{e} & \texttt{2} & \texttt{a} & \texttt{6} & \texttt{1} & \texttt{d} & \texttt{f} & \texttt{3} & \texttt{4} & \texttt{8}\\
\texttt{d} & \texttt{0} & \texttt{d} & \texttt{9} & \texttt{4} & \texttt{1} & \texttt{c} & \texttt{8} & \texttt{5} & \texttt{2} & \texttt{f} & \texttt{b} & \texttt{6} & \texttt{3} & \texttt{e} & \texttt{a} & \texttt{7}\\
\texttt{e} & \texttt{0} & \texttt{e} & \texttt{f} & \texttt{1} & \texttt{d} & \texttt{3} & \texttt{2} & \texttt{c} & \texttt{9} & \texttt{7} & \texttt{6} & \texttt{8} & \texttt{4} & \texttt{a} & \texttt{b} & \texttt{5}\\
\texttt{f} & \texttt{0} & \texttt{f} & \texttt{d} & \texttt{2} & \texttt{9} & \texttt{6} & \texttt{4} & \texttt{b} & \texttt{1} & \texttt{e} & \texttt{c} & \texttt{3} & \texttt{8} & \texttt{7} & \texttt{5} & \texttt{a}\\
\bottomrule
\end{tabular}
\caption{Producto $a\cdot b$ en $\FF4$, con $a$ en la fila y $b$ en la
columna (como en todas las tablas de este apéndice). Cada fila no
nula es una permutación de $\{1,\dots,\mathtt{f}\}$: la tabla lo
muestra y la estructura de cuerpo lo garantiza.}
\end{table}

\section{Cómo usar estas tablas}

\begin{itemize}[leftmargin=1.5em,itemsep=3pt]
  \item \textbf{Sumar} es XOR sin acarreo, bit a bit:
        $\mathtt{6}\oplus\mathtt{5}=\mathtt{3}$, es decir
        $(z+z^2)+(1+z^2)=1+z$.
  \item \textbf{Multiplicar} pasa por logaritmos:
        $a\,b=z^{\log_z a+\log_z b \pmod{15}}$. Ejemplo:
        $\mathtt{6}\cdot\mathtt{5}=z^{5}z^{8}=z^{13}=\mathtt{d}$;
        la versión polinómica es
        $(z+z^2)(1+z^2)=1+z^2+z^3$.
  \item \textbf{Invertir} es cambiar de signo el logaritmo. Con eso,
        la división es una resta de exponentes, y los
        $1/g(\alpha_i)$ del capítulo~\ref{ch:labgoppa} se calculan
        sin escribir una sola división polinómica.
  \item \textbf{La raíz cuadrada} no necesita algoritmos iterativos:
        en característica $2$, $\sqrt{x}=x^{2^{m-1}}=x^8$; Patterson
        la usa constantemente (capítulo~\ref{ch:goppa}).
  \item \textbf{La traza} decide irreducibilidad y pertenencia a
        subcuerpos; en el juguete, $\Tr(\beta)=\Tr(\beta^2)$ para
        todo $\beta$ (ejercicio~\ref{ej:traza-ej}).
  \item \textbf{Verificación cruzada.} El programa del
        apéndice~\ref{app:codigo} construye estas mismas tablas en
        una función de dos líneas y las usa para comprobar, número a
        número, todas las cuentas de los
        capítulos~\ref{ch:labgoppa} y~\ref{ch:labkem}.
\end{itemize}
